\section{Methodology}
\label{sec:method}

\subsection{Experimental design}

We adopt a within-subjects factorial design with two factors: preprocessing strategy (five levels) and retriever architecture (two levels). All queries are evaluated under all combinations of conditions, except Strategy~5 (Oracle), which is evaluated only on a manually annotated subset of approximately 150 queries.

\subsection{Hypotheses}

We pre-register five falsifiable hypotheses:

% =============================================================================
% shared/hypotheses.tex
% Single responsibility: canonical statement of the four pre-registered hypotheses.
% Change this file when a hypothesis text or test is updated.
% \input from any document that needs to state the hypotheses.
% =============================================================================

\noindent\textbf{H1 (omnibus).} At least one of the four preprocessing strategies yields a significantly different per-query nDCG@10 distribution than the others, separately for each retriever architecture. \textit{Test:} Friedman's test, $\alpha = 0.05$.

\smallskip
\noindent\textbf{H2 (harm of naive strip).} Naive regex \id{-nya} stripping yields significantly lower nDCG@10 than Keep on both retrievers, due to over-stripping of high-frequency non-clitic words such as \id{punya}, \id{tanya}, and \id{hanya}. \textit{Test:} one-tailed paired Wilcoxon signed-rank, $\alpha = 0.05$.

\smallskip
\noindent\textbf{H3 (architecture interaction).} The optimal preprocessing strategy among Keep, Sastrawi clitic, and Sentinel differs between BM25 and BGE-m3. \textit{Test:} comparison of best-performing strategy per retriever with bootstrap confidence intervals on the gain difference.

\smallskip
\noindent\textbf{H4 (sensitivity stratification).} The magnitude of differences between strategies is larger on high-\id{-nya}-sensitivity queries than on low-sensitivity queries. \textit{Test:} mixed-effects model with strategy-by-stratum interaction term.


\subsection{Preprocessing conditions}

Each strategy is applied identically to queries and passages.

% =============================================================================
% shared/preprocessing_strategies.tex
% Single responsibility: canonical enumeration of the 4 preprocessing conditions.
% Change this file when a strategy is added/removed or its specification changes.
% \input from any document that needs to enumerate the strategies.
% =============================================================================

\begin{enumerate}
    \item \textbf{Keep} --- identity preprocessing; UTF-8 normalised to NFC. Baseline.
    \item \textbf{Naive strip} --- regex \texttt{(\textbackslash w+?)nya\textbackslash b} replaced with the captured group. Included as a deliberate harm-baseline; over-strips \id{punya}, \id{tanya}, \id{hanya}, \id{biasanya}, and proper nouns ending in \id{-nya}.
    \item \textbf{Sastrawi clitic} --- Sastrawi's \texttt{RemovePossessivePronoun} filter applied with an explicit Indonesian-dictionary guard; strips \id{-nya} only when the result is a known root, protecting non-clitic words.
    \item \textbf{Sentinel} --- \id{-nya} replaced with a spaced sentinel token \texttt{<NYA>}. Diagnostic condition that removes lexical match while preserving the information that a clitic was present.
\end{enumerate}


Implementation details for each strategy --- including dictionary-guard wrapping for Strategy~3 and the false-positive verification protocol for Strategy~2 --- are reported in \secref{sec:experiments}.

\subsection{Retrievers}

We evaluate two retriever architectures spanning the sparse--dense spectrum:

\paragraph{BM25.} Standard BM25 implementation in pyserini, with $k_1 = 0.9$ and $b = 0.4$ (standard MIRACL configuration).

\paragraph{BGE-m3.} The dense embedding mode of BGE-m3 \citep{chen2024bgem3}, with maximum sequence length capped at 512 tokens. Indices built using FAISS HNSW (M=32, efConstruction=200).

\subsection{Dataset}

We evaluate on the Indonesian split of MIRACL \citep{zhang2023miracl}: approximately 1.4 million Wikipedia passages and 960 development-set queries with native-annotated relevance judgments.

\subsection{Metrics}

\textbf{Primary:} nDCG@10. \textbf{Secondary:} MRR@100, nDCG@1, Recall@10, Recall@100. nDCG@1 is reported as a more sensitive metric in case of saturation effects on dense baselines.

\subsection{Statistical analysis}

Per-query metrics are computed for all conditions. Omnibus tests use Friedman's test on per-query nDCG@10. Pairwise post-hoc tests use Wilcoxon signed-rank with Bonferroni correction across the $\binom{5}{2} = 10$ pairs per retriever ($\alpha = 0.005$ within each retriever family). Effect sizes are reported as Cliff's $\delta$. Bootstrap 95\% confidence intervals on $\Delta$nDCG@10 are computed with 10{,}000 resamples. The H5 stratification analysis uses a mixed-effects model with strategy-by-stratum interaction.

\subsection{Reproducibility and pre-registration}

This methodology was committed to a public Git repository before any experimental code was run, serving as a lightweight pre-registration of the design. All code, intermediate results, and analysis scripts will be released under an open licence on completion.
